% ============================================================================
%  B/09-nghe-giang — file chương
%  Ngày tạo: 2026-09-06 | Sửa gần nhất: 2026-09-06 | Ôn gần nhất: chưa
%  Dependency: B/06. Mức độ: cốt lõi cho người học qua khoá học trực tuyến.
% ============================================================================
\documentclass[../../english.tex]{subfiles}
\begin{document}
\chapter{Nghe giảng và xem khoá học (Following Lectures and Courses)}
\ptrtarget{B/09}\ptrtarget{B/09-nghe-giang}
\dependency{\ptr{B/06} cho tín hiệu cấu trúc --- bài giảng dùng cùng bộ tín hiệu đó nhưng ở dạng nói; \ptr{A/02} cho ngưỡng từ vựng.}

\begin{tomtat}
Nghe khác đọc ở ba điểm quyết định: bạn không dừng lại tra được, thông tin trôi theo thời gian, và một phần nghĩa nằm ở trọng âm và ngữ điệu chứ không ở từ. Chương này đưa ra bốn chiến lược đã được kiểm nghiệm cho việc theo một khoá học kỹ thuật bằng tiếng Anh: xem trước tài liệu để chuyển gánh nặng từ nghe sang đọc; nghe theo tín hiệu cấu trúc thay vì nghe từng từ; dùng phụ đề có phương pháp thay vì bật mặc định; và ghi chú theo khung chứ không chép lời. Chương cũng giải thích vì sao bạn không nghe ra những từ mà bạn biết rõ khi đọc --- hiện tượng nói nối và dạng yếu --- và làm gì với nó. Nếu chỉ nhớ một điều: mười phút đọc trước bài giảng tiết kiệm được nhiều hơn một giờ nghe lại.
\end{tomtat}

\begin{nentang}
\kn{tín hiệu cấu trúc}{các cụm mà người giảng dùng để báo chuyển ý: \emph{let's move on to}, \emph{the key point here is}, \emph{to summarise}, \emph{a common mistake is}.}
\kn{trọng âm câu}{trong tiếng Anh, từ mang thông tin mới được nhấn mạnh; vị trí trọng âm cho biết tác giả đang đối lập hay nhấn mạnh điều gì.}
\kn{nói nối}{trong lời nói liên tục, các từ dính vào nhau: \emph{an example} nghe như \emph{a-nex-ample}. Đây là lý do chính bạn không nhận ra từ quen.}
\kn{dạng yếu}{các từ chức năng bị rút gọn khi không nhấn: \emph{of} thành \emph{əv}, \emph{can} thành \emph{kən}. Người học thường chờ nghe dạng đầy đủ nên bỏ sót.}
\kn{ghi chú theo khung}{ghi dàn ý trước khi nghe (từ slide hoặc mục lục), rồi chỉ điền nội dung vào khung trong lúc nghe.}
\kn{xem trước}{đọc tài liệu, slide hoặc bài báo tương ứng \emph{trước} khi nghe, để chuyển phần lớn gánh nặng hiểu từ kênh nghe sang kênh đọc.}
\end{nentang}

\begin{khoigoi}
\h{Vì sao bạn nghe không ra một từ mà nếu nhìn thấy viết ra thì bạn biết ngay?}
\h{Tại sao mười phút đọc trước lại hiệu quả hơn ba mươi phút nghe lại?}
\h{Bật phụ đề có làm bạn học kém đi không, và nếu có thì nên dùng thế nào?}
\h{Ghi chú trong lúc nghe làm bạn bỏ lỡ nội dung. Vậy nên ghi thế nào?}
\h{Trọng âm trong câu tiếng Anh mang thông tin gì mà chữ viết không có?}
\end{khoigoi}

Có một trải nghiệm rất quen với người học tiếng Anh khá: bạn đọc một chương sách kỹ thuật không mấy khó khăn, nhưng xem bài giảng về đúng chủ đề đó thì mất mạch sau mười phút. Cảm giác là ``kỹ năng nghe của tôi yếu''.

Chẩn đoán đó thường đúng một nửa. Vấn đề không hẳn ở khả năng nhận diện âm thanh mà ở \emph{ngân sách chú ý}. Khi đọc, bạn kiểm soát tốc độ: gặp chỗ khó thì chậm lại, quên thì nhìn lại đoạn trước. Khi nghe, tốc độ do người khác quyết định, và nếu bạn dừng lại xử lý một câu thì hai câu tiếp theo trôi mất. Vì vậy chiến lược cho việc nghe không phải ``nghe kỹ hơn'' mà là \emph{giảm tải}: làm sao để lúc nghe bạn còn dư chú ý.

Bốn chiến lược trong chương này đều là các cách giảm tải, và cái hiệu quả nhất --- xem trước --- lại là cái ít người làm nhất.

\section{Vì sao bạn không nghe ra từ mình đã biết}

Câu hỏi mở đầu thứ nhất có một câu trả lời khá kỹ thuật, và biết nó giúp bạn ngừng tự trách.

Trong lời nói liên tục, các từ \emph{không} được phát âm như trong từ điển. Ba hiện tượng chính.

\emph{Nói nối.} Âm cuối của từ này dính vào âm đầu của từ sau: \emph{look at it} nghe thành một khối \emph{lookatit}. Bạn đang chờ ba từ tách rời và nghe thấy một chuỗi âm lạ.

\emph{Dạng yếu.} Các từ chức năng --- giới từ, mạo từ, trợ động từ --- bị rút gọn tới mức gần như biến mất khi không mang trọng âm. \emph{We can do it} thường nghe như \emph{we kən do it}, và người học chờ nghe rõ \emph{can} nên bỏ sót cả câu. Đáng chú ý: chính vì \emph{can} bị rút gọn còn \emph{can't} thì không, nhiều người nghe nhầm hai câu trái nghĩa nhau.

\emph{Đồng hoá và lược âm.} Âm biến đổi hoặc rơi hẳn: \emph{next day} thành \emph{neks-day}, \emph{going to} thành \emph{gonna} trong lời nói không trang trọng.

Hệ quả thực hành: luyện nghe không phải luyện ``nghe rõ hơn'' mà là làm quen với \emph{dạng thật} của các từ trong dòng lời nói. Cách hiệu quả nhất là đọc theo bản ghi: nghe một câu, nhìn bản ghi, nghe lại, và chú ý xem chỗ nào khác với hình dung của bạn. Vài chục câu như vậy đủ để nhận ra các mẫu lặp lại, vì các hiện tượng trên rất có quy luật.

\section{Chiến lược 1: xem trước --- đòn bẩy lớn nhất}

Đây là câu trả lời cho câu hỏi mở đầu thứ hai và là lời khuyên có tỉ lệ lợi ích trên công sức cao nhất chương.

Khi nghe một bài giảng kỹ thuật, bạn phải làm ba việc cùng lúc: nhận diện âm thanh, hiểu nội dung chuyên môn, và theo dõi cấu trúc lập luận. Ba việc đó tranh nhau một ngân sách chú ý có hạn. Nếu bạn đã biết trước nội dung chuyên môn, hai phần ba gánh nặng biến mất và phần nghe trở nên dễ hơn hẳn --- không phải vì tai bạn tốt lên mà vì bạn \emph{dự đoán} được điều sắp nghe.

Cách làm cụ thể, mười phút cho một bài giảng một giờ: đọc lướt slide nếu có; hoặc đọc tóm tắt và các đề mục của bài báo tương ứng; hoặc đọc mục lục chương sách. Ghi ra ba câu hỏi bạn muốn bài giảng trả lời. Ba câu hỏi đó biến việc nghe từ tiếp nhận thụ động thành tìm kiếm có mục tiêu, và đó là khác biệt lớn nhất.

Với các khoá học có phụ đề hoặc bản ghi, còn một biến thể mạnh hơn: đọc bản ghi trước ở tốc độ của bạn, rồi nghe. Lần nghe khi đó gần như chỉ còn là củng cố --- và nó cũng chính là cách luyện nhận diện âm thanh ở mục 1.

\section{Chiến lược 2: nghe cấu trúc, không nghe từ}

Người giảng có kinh nghiệm rắc tín hiệu cấu trúc dày đặc, và các tín hiệu này dễ nghe hơn nội dung vì chúng là những cụm cố định lặp lại.

\begin{center}\footnotesize
\rowcolors{2}{white}{tblalt}
\begin{longtable}{@{}L{3.6cm}L{5.4cm}L{5.4cm}@{}}
\toprule
\textbf{Chức năng} & \textbf{Cụm hay dùng} & \textbf{Bạn nên làm gì khi nghe thấy}\\
\midrule
\endfirsthead
\toprule
\textbf{Chức năng} & \textbf{Cụm hay dùng} & \textbf{Bạn nên làm gì (tiếp)}\\
\midrule
\endhead
Báo dàn ý & \emph{Today I'll cover three things} & Ghi ba đầu mục làm khung ghi chú\\
Chuyển ý & \emph{Let's move on to}, \emph{Turning to} & Xuống dòng mới trong ghi chú\\
Nhấn trọng tâm & \emph{The key point here is}, \emph{If you remember one thing} & Ghi nguyên văn --- đây là câu đáng nhớ nhất\\
Định nghĩa & \emph{What we mean by X is}, \emph{X is defined as} & Ghi định nghĩa vào bảng thuật ngữ (\ptr{B/10})\\
Ví dụ & \emph{For instance}, \emph{Let me give you an example} & Có thể nghe lơi; ví dụ thường minh hoạ ý vừa nói\\
Cảnh báo lỗi & \emph{A common mistake is}, \emph{People often think that} & Ghi lại --- đây thường là phần giá trị nhất của bài giảng\\
Ngoài lề & \emph{As an aside}, \emph{This is not on the exam} & Được phép thư giãn chú ý\\
Tóm tắt & \emph{To summarise}, \emph{So the takeaway is} & Đối chiếu với khung ghi chú của bạn\\
\bottomrule
\end{longtable}
\end{center}

Nguyên tắc: khi mất mạch, đừng cố dịch câu vừa nghe --- hãy \emph{chờ tín hiệu cấu trúc tiếp theo} và bắt lại từ đó. Bài giảng có độ dư thừa cao hơn văn viết rất nhiều; điều quan trọng gần như luôn được nói lại.

\section{Chiến lược 3: dùng phụ đề có phương pháp}

Câu hỏi mở đầu thứ ba là một câu hỏi thật. Bật phụ đề giúp bạn hiểu ngay, nhưng nếu luôn bật thì bạn đang \emph{đọc} chứ không luyện nghe, và khả năng nhận diện âm thanh không tiến.

Ba cách dùng có phương pháp, chọn theo mục tiêu của buổi học.

\emph{Khi mục tiêu là nắm nội dung chuyên môn:} bật phụ đề, không áy náy. Đừng biến buổi học kỹ thuật thành buổi luyện nghe --- hai mục tiêu tranh nhau và bạn hỏng cả hai.

\emph{Khi mục tiêu là luyện nghe:} nghe một đoạn không phụ đề; nghe lại lần hai; rồi bật phụ đề và đối chiếu; rồi nghe lần cuối không phụ đề. Đoạn nên ngắn --- hai tới ba phút. Chính bước đối chiếu là bước dạy bạn nhiều nhất, vì nó chỉ ra chính xác chỗ bạn nghe nhầm.

\emph{Khi tài liệu quá khó:} bật phụ đề và giảm tốc độ xuống 0,75. Nhưng đừng ở lâu tại đó --- tốc độ chậm tạo thói quen xử lý chậm. Khi đã quen, hãy tăng dần, và nhiều người thấy 1,25 lần lại tốt hơn cho việc giữ tập trung.

\section{Chiến lược 4: ghi chú theo khung}

Câu hỏi mở đầu thứ tư nêu một mâu thuẫn thật: ghi chú giúp nhớ nhưng lại chiếm chú ý đang cần cho việc nghe. Lời giải là \emph{chuẩn bị khung trước}, để lúc nghe chỉ phải điền.

Quy trình ba giai đoạn.

\emph{Trước:} dựng khung từ slide hoặc đề mục --- các tiêu đề, mỗi tiêu đề chừa vài dòng trống. Nếu không có slide, hãy để trống và tạo mục mới mỗi khi nghe tín hiệu chuyển ý.

\emph{Trong:} ghi \emph{ít}. Ghi các định nghĩa, các con số, các cảnh báo lỗi, và những câu người giảng nhấn mạnh. Không chép lại những gì đã có trên slide --- đó là lãng phí chú ý lớn nhất. Dùng lề trái cho câu hỏi nảy ra trong đầu.

\emph{Sau:} dành năm phút ngay sau buổi học viết một đoạn tóm tắt bằng lời của bạn, và trả lời ba câu hỏi bạn đặt ra ở bước xem trước. Năm phút này quyết định phần lớn việc bạn còn nhớ gì sau một tuần --- vì nó là một lần \emph{truy xuất} chủ động chứ không phải đọc lại (\ptr{A/02}).

Bố cục hai cột --- ghi chú bên phải, câu hỏi và từ khoá bên trái, tóm tắt ở dưới cùng --- là một khuôn phổ biến và hiệu quả, vì cột trái sau này dùng làm mặt trước của thẻ ôn tập.

\begin{cannho}
Ba việc quyết định hiệu quả của một buổi học qua video, và cả ba nằm ngoài lúc xem: mười phút xem trước tài liệu và viết ba câu hỏi; dựng khung ghi chú trước khi bấm phát; và năm phút sau buổi học viết tóm tắt bằng lời của mình. Nếu chỉ làm được một, hãy làm việc đầu tiên.
\end{cannho}

\section{Trọng âm mang nghĩa}

Câu hỏi mở đầu thứ năm chạm vào một tầng thông tin mà chữ viết không có. Trong tiếng Anh, từ mang \emph{thông tin mới} được nhấn mạnh, và người nghe dùng vị trí trọng âm để biết trọng tâm câu nằm ở đâu.

So sánh cùng một câu với ba vị trí nhấn khác nhau --- phần in đậm là phần được nhấn:

\emph{We tested it on \textbf{three} datasets} --- nhấn số lượng, ngụ ý ba là điểm đáng chú ý (có thể là ít, có thể là nhiều).

\emph{We \textbf{tested} it on three datasets} --- nhấn hành động, ngụ ý đối lập với việc chỉ phân tích lý thuyết.

\emph{\textbf{We} tested it on three datasets} --- nhấn chủ thể, ngụ ý đối lập với người khác: chúng tôi làm, không phải họ.

Ba câu giống hệt nhau trên giấy, ba nghĩa khác nhau khi nói. Đây cũng là lý do nghe hiểu không thể quy về đọc: một phần nội dung chỉ tồn tại ở kênh âm thanh.

Liên hệ đáng chú ý với phần viết: quy tắc ``thông tin mới được nhấn'' trong lời nói tương ứng với quy tắc ``đặt thông tin mới ở cuối câu'' trong văn viết --- vì cuối câu là vị trí nhấn tự nhiên của văn viết. Chương \ptr{C/12} khai thác trọn vẹn quy tắc đó.

\section{Gốc rễ: từ ghi chép giảng đường tới khoá học trực tuyến}

Việc học qua nghe giảng cổ hơn nhiều so với việc học qua đọc: trước khi in ấn phổ biến, giảng đường là nơi tri thức được truyền, và \emph{lecture} trong tiếng Anh có gốc Latin nghĩa là ``đọc'' --- thầy đọc bản thảo cho sinh viên chép, vì mỗi bản sao đều phải chép tay.

Điều đó để lại dấu vết trong cấu trúc bài giảng tới hôm nay: độ dư thừa cao, lặp lại có chủ ý, và các tín hiệu chuyển ý rõ ràng --- tất cả đều là đặc điểm của ngôn ngữ nói được thiết kế để nghe một lần. Khi đọc, bạn có thể nhìn lại; khi nghe, người nói phải bù bằng cách lặp.

Bố cục ghi chú hai cột được đề xuất tại Đại học Cornell vào thập niên 1950, trong bối cảnh giáo dục đại học Mỹ mở rộng nhanh và sinh viên cần một phương pháp học có hệ thống. Ý tưởng cốt lõi không phải bố cục mà là \emph{giai đoạn sau buổi học}: cột trái được điền sau, bằng các câu hỏi rút ra từ ghi chú, và chính việc trả lời các câu hỏi đó mới là học. Nói theo ngôn ngữ của \ptr{A/02}: bố cục ép bạn truy xuất chủ động thay vì đọc lại.

Bước ngoặt gần đây là khoá học trực tuyến. Nó đổi hai điều theo hướng có lợi cho người học ngoại ngữ: bạn kiểm soát được tốc độ và có thể nghe lại, nên rào cản âm thanh giảm hẳn; và bản ghi cùng phụ đề cho phép kết hợp kênh nghe với kênh đọc theo cách mà giảng đường không cho. Nhưng nó cũng lấy đi một thứ: không có ai để hỏi lại khi không hiểu, nên phần \emph{sau buổi học} --- tóm tắt, ghi câu hỏi, tìm câu trả lời --- trở nên quan trọng hơn nhiều so với thời chỉ có giảng đường.

\begin{gocre}
\gr{Thời trung cổ --- \emph{lecture} nghĩa gốc là ``đọc''}{Sách chép tay quá đắt; tri thức truyền qua giọng nói.}{Bài giảng hình thành với độ dư thừa cao và tín hiệu chuyển ý rõ --- đặc điểm của ngôn ngữ để nghe một lần.}
\gr{Thập niên 1950 --- bố cục ghi chú hai cột}{Sinh viên chép lại mọi thứ mà không nhớ được gì.}{Tách ghi chú khỏi \emph{câu hỏi} và \emph{tóm tắt}; ép truy xuất chủ động sau buổi học --- cơ chế được nghiên cứu về hiệu ứng kiểm tra xác nhận\cite{roediger2006}.}
\gr{Nghiên cứu về diễn ngôn bài giảng\cite{biber1999}}{Người học ngoại ngữ nghe không kịp dù đủ từ vựng.}{Chỉ ra vai trò của tín hiệu cấu trúc và của độ dư thừa; gợi ý dạy nghe theo cấu trúc thay vì theo từ.}
\gr{Thập niên 2010 --- khoá học trực tuyến}{Người học không tiếp cận được giảng đường tốt.}{Kiểm soát tốc độ, nghe lại, phụ đề và bản ghi --- rào cản âm thanh giảm, nhưng mất kênh hỏi lại nên khâu sau buổi học quan trọng hơn.}
\end{gocre}

\section{Liên hệ ngang}

\begin{lienhe}
\lh{Âm nhạc}{Luyện tai: nghe hợp âm, nhận ra tiết tấu}{Cùng cơ chế: bạn chỉ nghe ra thứ mà bạn đã có mô hình trong đầu. Đó là lý do xem trước hiệu quả tới vậy --- nó dựng sẵn mô hình để tai khớp vào.}
\lh{Họp hành công việc}{Theo dõi cuộc họp bằng tiếng Anh}{Cùng bộ tín hiệu cấu trúc và cùng chiến lược: chuẩn bị trước bằng cách đọc chương trình họp, và bắt lại từ tín hiệu chuyển ý khi mất mạch. Khác: họp có thể hỏi lại, nên hãy hỏi.}
\lh{Đọc mã cùng người khác}{Nghe đồng nghiệp giải thích một đoạn mã}{Người giải thích cũng dùng tín hiệu chuyển ý và nhấn mạnh. Mẹo tương tự: xem mã trước khi nghe giải thích.}
\lh{Ngôn ngữ học âm vị}{Nói nối, đồng hoá, dạng yếu}{Các hiện tượng ở mục 1 có quy luật rõ, không phải ``người bản ngữ nói ẩu''. Biết quy luật giúp bạn dự đoán được dạng thật của từ trong dòng lời nói.}
\lh{Nhận dạng tiếng nói}{Mô hình ngôn ngữ giúp nhận dạng âm thanh}{Hệ thống nhận dạng dùng kiến thức về từ nào có thể theo sau từ nào để bù cho tín hiệu mờ --- đúng việc mà bạn làm khi xem trước tài liệu. Cùng nguyên lý: dự đoán bù cho tín hiệu kém.}
\lh{Giảng dạy}{Thiết kế bài giảng có tín hiệu rõ}{Khi bạn đứng ở phía người giảng, danh sách ở mục 3 trở thành danh sách những cụm bạn \emph{nên} dùng --- đặc biệt khi người nghe không dùng tiếng Anh bản ngữ.}
\end{lienhe}

\begin{traloi}
\tl{Vì trong lời nói liên tục, từ không được phát âm như trong từ điển: âm cuối dính vào từ sau (nói nối), các từ chức năng bị rút gọn thành dạng yếu, và nhiều âm bị đồng hoá hoặc rơi hẳn. Bạn đang chờ dạng đầy đủ và nghe thấy một chuỗi âm khác. Cách chữa không phải ``nghe kỹ hơn'' mà là làm quen với dạng thật: nghe một câu, đối chiếu bản ghi, nghe lại, và chú ý chỗ khác với hình dung của bạn. Các hiện tượng này có quy luật nên vài chục câu là đủ để bắt được mẫu.}
\tl{Vì ba việc --- nhận diện âm, hiểu nội dung chuyên môn, theo dõi cấu trúc --- tranh nhau một ngân sách chú ý có hạn. Đọc trước xoá phần lớn gánh nặng thứ hai, và quan trọng hơn, nó cho bạn khả năng \emph{dự đoán} câu tiếp theo, khiến việc nhận diện âm dễ hẳn. Nghe lại thì không giảm tải: bạn vẫn phải làm cả ba việc cùng lúc, chỉ là làm hai lần.}
\tl{Có, nếu luôn bật --- vì khi đó bạn đang đọc chứ không luyện nghe. Cách dùng theo mục tiêu: khi cần nắm nội dung chuyên môn thì cứ bật, đừng biến buổi học kỹ thuật thành buổi luyện nghe; khi muốn luyện nghe thì nghe đoạn ngắn không phụ đề hai lần, rồi bật phụ đề đối chiếu, rồi nghe lại không phụ đề --- bước đối chiếu là bước dạy nhiều nhất; và khi tài liệu quá khó thì bật kèm giảm tốc, nhưng chỉ tạm thời.}
\tl{Chuẩn bị khung trước khi nghe, để trong lúc nghe chỉ phải điền. Ghi ít: định nghĩa, con số, cảnh báo lỗi, và những câu người giảng nhấn mạnh --- tuyệt đối không chép lại những gì đã có trên slide. Dành lề trái cho câu hỏi nảy ra. Rồi năm phút sau buổi học, viết một đoạn tóm tắt bằng lời của mình và trả lời các câu hỏi đã đặt ra --- đây là phần quyết định bạn còn nhớ gì sau một tuần.}
\tl{Trọng âm đánh dấu \emph{thông tin mới} hoặc \emph{đối lập}. Cùng một câu viết ra, nhấn vào ``ba'' thì trọng tâm là số lượng; nhấn vào động từ thì đối lập với việc chỉ làm lý thuyết; nhấn vào chủ ngữ thì đối lập với người khác. Chữ viết không mang tầng này, và đó là lý do nghe hiểu không quy về đọc được. Liên hệ với viết: quy tắc ``mới thì nhấn'' trong nói tương ứng với quy tắc ``đặt thông tin mới ở cuối câu'' trong viết (\ptr{C/12}).}
\end{traloi}

\begin{dongian}
Nghe khó hơn đọc không phải vì tai bạn kém, mà vì bạn không được quyền dừng lại. Khi đọc, gặp chỗ khó thì chậm lại; khi nghe, dừng lại nghĩ một câu là mất hai câu sau.

Vì vậy mẹo quan trọng nhất không nằm ở lúc nghe mà ở trước đó: hãy dành mười phút đọc slide hoặc đọc tóm tắt bài báo tương ứng, rồi viết ra ba câu hỏi bạn muốn được trả lời. Khi đã biết trước nội dung sẽ nói về gì, tai bạn chỉ còn việc nhận ra từ --- và nó làm việc đó tốt hơn nhiều khi có thể đoán trước.

Mẹo thứ hai: đừng cố nghe từng từ. Người giảng luôn rắc các câu chỉ đường --- \emph{today I'll cover three things}, \emph{the key point here is}, \emph{a common mistake is}. Những câu này dễ nghe vì chúng lặp đi lặp lại, và chúng cho bạn biết mình đang ở đâu. Khi mất mạch, đừng dịch lại câu vừa rồi; hãy chờ câu chỉ đường tiếp theo rồi bắt lại. Bài giảng nói lại điều quan trọng nhiều lần, nên bạn sẽ không mất nó.

Còn chuyện vì sao bạn không nghe ra những từ mình biết rõ: vì trong lời nói, các từ dính vào nhau và các từ nhỏ bị nuốt. \emph{Look at it} nghe thành một khối. \emph{Can} bị rút gọn gần như biến mất, trong khi \emph{can't} thì không --- nên hai câu ngược nghĩa lại dễ nghe nhầm. Đây là quy luật chứ không phải người ta nói ẩu, và cách quen với nó là nghe kèm bản ghi rồi đối chiếu.

Cuối cùng, năm phút sau buổi học đáng giá hơn ba mươi phút xem lại: gấp máy lại và viết một đoạn tóm tắt bằng lời của mình.
\motcau{Mười phút đọc trước tiết kiệm hơn một giờ nghe lại --- vì nghe dễ khi bạn đã đoán được điều sắp nghe.}
\chuadongian{Giọng vùng miền lạ và chất lượng âm thanh kém là rào cản thật; với chúng, không mẹo nào thay được việc nghe nhiều giờ để quen.}
\end{dongian}

\begin{onbai}
\ar[3]{Ba khác biệt giữa nghe và đọc}{Không dừng lại được; thông tin trôi theo thời gian; một phần nghĩa nằm ở trọng âm và ngữ điệu.}
\ar[3]{Chiến lược mạnh nhất}{Xem trước: đọc slide hoặc tóm tắt và viết ba câu hỏi --- giảm tải chú ý và cho khả năng dự đoán.}
\ar[3]{Ba hiện tượng làm bạn nghe không ra từ quen}{Nói nối (từ dính nhau); dạng yếu (từ chức năng bị rút gọn); đồng hoá và lược âm. Đều có quy luật, học được.}
\ar[3]{Nghe cấu trúc, không nghe từ}{Bắt các tín hiệu: \emph{today I'll cover}, \emph{let's move on}, \emph{the key point is}, \emph{a common mistake is}. Mất mạch thì chờ tín hiệu tiếp theo.}
\ar[3]{Ba cách dùng phụ đề}{Nắm nội dung: cứ bật. Luyện nghe: nghe không phụ đề hai lần, đối chiếu, rồi nghe lại. Quá khó: bật kèm giảm tốc, nhưng tạm thời.}
\ar[3]{Ghi chú ba giai đoạn}{Trước: dựng khung từ slide. Trong: ghi ít, không chép slide. Sau: năm phút tóm tắt bằng lời mình --- phần quyết định trí nhớ.}
\ar[2]{Trọng âm mang nghĩa}{Từ mang thông tin mới được nhấn; vị trí nhấn cho biết trọng tâm và đối lập --- tầng thông tin mà chữ viết không có.}
\ar[2]{Liên hệ nói--viết}{``Mới thì nhấn'' trong nói tương ứng ``mới thì đặt cuối câu'' trong viết (\ptr{C/12}).}
\ar[2]{\emph{Can} và \emph{can't}}{\emph{Can} bị rút gọn mạnh, \emph{can't} thì không --- nguồn nghe nhầm hai câu trái nghĩa.}
\ar[2]{Vì sao bài giảng dư thừa}{Ngôn ngữ nói được thiết kế để nghe một lần, nên người nói bù bằng cách lặp --- bạn được phép bỏ lỡ và bắt lại.}
\ar[1]{Tốc độ phát}{0,75 khi quá khó nhưng đừng ở lâu; nhiều người thấy 1,25 giúp giữ tập trung tốt hơn khi đã quen.}
\end{onbai}

\begin{hoidap}
\qa{Tôi nên xem khoá học bằng tiếng Anh hay tìm bản tiếng Việt?}{Nếu mục tiêu là nắm khái niệm mới hoàn toàn, xem tiếng Việt trước rồi xem lại tiếng Anh là hợp lý --- lần hai bạn đã có khung khái niệm nên chỉ còn học ngôn ngữ. Nếu chủ đề bạn đã có nền, hãy xem thẳng tiếng Anh, vì đó là nơi thuật ngữ và cách diễn đạt của ngành tồn tại. Về lâu dài, mọi tài liệu chuyên sâu đều bằng tiếng Anh, nên chuyển sang càng sớm càng tốt --- với điều kiện dùng chiến lược xem trước để không bị quá tải.}
\qa{Tôi nghe hiểu nhưng không nhớ được gì sau buổi học. Vấn đề ở đâu?}{Ở khâu sau buổi học, gần như chắc chắn. Nghe hiểu là xử lý tức thời và nó không tự chuyển thành trí nhớ dài hạn; cần một lần truy xuất chủ động. Năm phút viết tóm tắt bằng lời của bạn --- không nhìn ghi chú --- là can thiệp rẻ nhất và hiệu quả nhất. Nếu vẫn không đủ, hãy thêm một bước: hôm sau, nhìn cột câu hỏi trong ghi chú và tự trả lời trước khi đọc phần nội dung.}
\qa{Có nên luyện nói theo (shadowing) không, khi mục tiêu của tôi là đọc và viết?}{Có, ở mức nhỏ và vì một lý do gián tiếp: nói theo buộc bạn xử lý nói nối và dạng yếu bằng chính miệng mình, và điều đó cải thiện khả năng \emph{nhận ra} chúng khi nghe. Mười phút mỗi ngày với một đoạn ngắn là đủ. Đừng đặt mục tiêu phát âm chuẩn --- với mục tiêu đọc và viết, lợi ích chính của nói theo là ở tai chứ không ở miệng.}
\qa{Người giảng nói quá nhanh, tôi có nên bỏ qua khoá đó không?}{Trước khi bỏ, thử ba việc theo thứ tự. Xem trước tài liệu --- phần lớn cảm giác ``quá nhanh'' thực ra là quá tải nội dung chứ không phải tốc độ. Giảm tốc xuống 0,85 và bật phụ đề trong vài buổi đầu. Và xem lại chỉ những đoạn bạn đã đánh dấu là quan trọng, thay vì xem lại cả bài. Nếu sau ba việc đó vẫn không theo được, khoá học có thể đang giả định một nền kiến thức bạn chưa có --- vấn đề nằm ở nội dung, không ở tiếng Anh.}
\qa{Ghi chú bằng tiếng Anh hay tiếng Việt?}{Ghi thuật ngữ và các cụm cố định bằng tiếng Anh --- luôn luôn, vì đó là dạng bạn sẽ gặp lại và cần dùng. Phần diễn giải và suy nghĩ của bạn thì ghi bằng ngôn ngữ nào nhanh hơn, thường là tiếng Việt. Điều cần tránh là dịch thuật ngữ sang tiếng Việt rồi chỉ ghi bản dịch: bạn sẽ mất liên kết với dạng tiếng Anh và gặp lại vấn đề ở \ptr{B/10} về bản dịch không khớp hoàn toàn.}
\qa{Phần hỏi đáp cuối buổi giảng tôi gần như không theo được. Bình thường không?}{Rất bình thường, và có lý do khách quan: câu hỏi từ khán giả thường không được ghi âm rõ, người hỏi không chuẩn bị nên nói lộn xộn, và không có tín hiệu cấu trúc nào. Chiến lược thực dụng: nghe câu \emph{trả lời} thay vì câu hỏi --- người giảng giỏi hầu như luôn nhắc lại câu hỏi trước khi trả lời, và phần trả lời được tổ chức tốt hơn nhiều. Nếu không nhắc lại thì phần đó thường không đáng để bạn tốn công.}
\end{hoidap}

\begin{baitap}
\bt[1]{Xem trước tài liệu 10 phút và viết ba câu hỏi, rồi xem một bài giảng}{Khoá học bạn đang theo}{So sánh cảm giác với buổi học không xem trước.}
\bt[1]{Ghi lại mọi tín hiệu cấu trúc bạn nghe được trong một bài giảng}{Khoá học bất kỳ}{Đối chiếu với bảng ở mục 3; bổ sung các cụm riêng của người giảng đó.}
\bt[2]{Luyện đối chiếu: nghe đoạn 2 phút không phụ đề, rồi đọc bản ghi}{Video có phụ đề}{Ghi lại năm chỗ bạn nghe nhầm và phân loại: nói nối, dạng yếu, hay từ lạ.}
\bt[2]{Dùng bố cục ghi chú hai cột cho một buổi học}{Khoá học bạn đang theo}{Cột trái điền \emph{sau} buổi học, bằng câu hỏi rút từ ghi chú.}
\bt[2]{Viết tóm tắt 5 phút sau buổi học, không nhìn ghi chú}{Khoá học bất kỳ}{Sau đó so với ghi chú và đánh dấu chỗ bạn quên --- đó là chỗ cần ôn.}
\bt[3]{Luyện nói theo 10 phút mỗi ngày trong một tuần với một đoạn ngắn}{Video bất kỳ}{Mục tiêu là quen với nói nối và dạng yếu, không phải phát âm chuẩn.}
\bt[3]{Xem một bài giảng ở 1,25 tốc độ sau khi đã xem trước tài liệu}{Khoá học bạn đang theo}{Kiểm xem việc xem trước có bù được cho tốc độ cao hơn không.}
\end{baitap}

\section*{Kết chương và đọc tiếp}

Chương này chuyển việc nghe từ ``cố hiểu từng từ'' sang \emph{quản lý ngân sách chú ý}. Ba điều đáng mang theo: xem trước là đòn bẩy lớn nhất; bắt tín hiệu cấu trúc thay vì bắt từ; và năm phút sau buổi học quyết định bạn còn nhớ gì.

Chương cuối của phần B quay lại một rào cản đã xuất hiện ở mọi chương: thuật ngữ. Khi bạn đọc tài liệu ngành hoặc nghe giảng chuyên môn, phần lớn chỗ khựng lại không phải từ tiếng Anh khó mà là thuật ngữ mang nghĩa quy ước hẹp --- và nhiều thuật ngữ trong số đó là những từ thông thường được gán nghĩa mới. Chương \ptr{B/10} nói cách xây vốn từ đó một cách có hệ thống.
\end{document}
